\documentclass[9pt]{beamer}

\usecolortheme[RGB={55,110,144}]{structure} % prima c'era questo bluetto: 20,104,204  % blu elettrico: 20,10,244

\usepackage{bbm,color,xcolor}
\usepackage[utf8]{inputenc}

\usepackage{fontenc,fancyvrb}

\usepackage{amsmath}
\usepackage{amssymb}
\usepackage{amsfonts}
\usepackage{graphicx}
\usepackage{mathrsfs}
\usepackage{pgfplots}

\definecolor{mycolor}{rgb}{0.28,0.69,0.7}
\definecolor{mycolor1}{rgb}{0.7,0.75,0.71}
\definecolor{mycolor2}{rgb}{0.58,1,1}

\usepackage{pifont}% http://ctan.org/pkg/pifont
\newcommand{\cmark}{\ding{51}}%
\newcommand{\xmark}{\ding{55}}%

\usepackage{default,amsmath,comment,amssymb,amsthm,amscd,latexsym,tcolorbox,array}
\usepackage{amsfonts,mathrsfs,nomencl,hyperref,threeparttable,textcomp}
\usepackage[english]{babel}
\usepackage{graphicx,tabularx,tikz,setspace,subfigure,epigraph,lscape,color,booktabs,multirow}
\usepackage{fancybox,changes}
\usetikzlibrary{plotmarks}
\usetikzlibrary{matrix}
\usetikzlibrary{positioning}
\usetikzlibrary{mindmap,backgrounds}

%% strange style
\usepackage{times}
%\usepackage{mathptmx}
%\usepackage[scaled=.90]{helvet}
%\usepackage{courier}

\usepackage[labelformat=empty,labelsep=none]{caption}
%\renewcommand{\tablename}{}

%\def\softO{\ensuremath{{\mathcal{O}}{\,\tilde{ }\,}}}

\addtobeamertemplate{footline}{\hfill\insertframenumber/\inserttotalframenumber}

\usepackage{color,xcolor,colortbl}

\definecolor{mycolor}{rgb}{0.28,0.69,0.7}
\definecolor{mycolor1}{rgb}{0.7,0.75,0.71}
\definecolor{mycolor2}{rgb}{0.58,1,1}

%\normalfont

\renewcommand{\familydefault}{\rmdefault}

\newcommand{\red}[1]{{\bf \color{red} #1}}

\newcommand{\calD}{\mathcal{D}}
\newcommand{\calZ}{\mathcal{Z}}
\newcommand{\minrank}{{r}(A)}

\newcommand{\rrY}{{\color{red} \bf Y}}

\newcommand{\bc}{\begin{center}}
\newcommand{\ec}{\end{center}}
\newcommand{\cit}{\color{olive}\bf}


\setbeamercolor{upcol}{fg=black,bg=green}      %\setbeamercolor{upcol}{fg=black,bg=orange} prima c'era questo
\setbeamercolor{lowcol}{fg=black,bg=green!40}  %\setbeamercolor{lowcol}{fg=black,bg=orange!40}

\setbeamertemplate{navigation symbols}{}

\theoremstyle{definition}
\newtheorem{prob}{Problem}[section]

\theoremstyle{definition}
\newtheorem{prop}{Proposition}[section]

\theoremstyle{definition}
\newtheorem{teo}{Theorem}[section]

\theoremstyle{definition}
\newtheorem{defin}{Definition}[section]

\theoremstyle{definition}
\newtheorem{coroll}{Corollary}[section]

\theoremstyle{definition}
\newtheorem{oss}{Remark}[section]


\title[]{\bf SPECTRA : a Maple library for linear matrix inequalities \smallskip \smallskip \smallskip}
\author{{\bf Simone Naldi} \\ \vspace{0.1cm} Technische Universit\"at Dortmund}
\date{Goethe Universit\"at Frankfurt \\ \vspace{0.1cm} December 20, 2016}


\usetheme{default}

% MACROS

%% Fields - Rings
\def\QQ{{\mathbbmss{Q}}} % rational numbers
\def\RR{{\mathbbmss{R}}} % real numbers
\def\CC{{\mathbbmss{C}}} % complex numbers
\def\ZZ{{\mathbbmss{Z}}} % integer numbers
\def\NN{{\mathbbmss{N}}} % natural numbers

\def\bigO{{\mathcal{O}}}

%% Polynomial families
\def\f{{\mathbf{f}}}
\def\g{{\mathbf{g}}}
\def\h{{\mathbf{h}}}
\def\q{{\mathbf{q}}}

%% For the section FERMETURE.TEX + NOETHER.TEX + BOUNDARY.TEX
\def\fermF{{\mathcal{F}}} % families F_i of varieties of dimension <= i
\def\fermD{{D}} % the equidimensional variety of all section about the properness of projections
\def\fermS{{\mathcal{S}}} % subvarieties of the singular locus of \fermD
\def\fermW{{\mathcal{W}}} % polar varieties of \fermD
\def\fermI{{\mathcal{I}}} % families of ideals
\def\fermT{{\mathcal{T}}} % a generic variety in F_i
\def\fermL{{\mathcal{L}}} % a set in the proof of boundary.tex
\def\fermU{{\mathfrak{A}}} % the parameters for the matrix \bfM in the proof of noether position
\def\fermP{{\mathcal{P}}} % a prime ideal
\def\fermQ{{\mathcal{Q}}} % a prime ideal
\def\sing{{\sf Sing}}
\def\crit{{\sf Crit}}
\def\reg{{\sf Reg}}
\def\scC{{\mathscr{C}}}
\def\scS{{\mathscr{S}}}
\def\sfP{{\mathsf{P}}}
\def\sfQ{{\mathsf{Q}}}
\def\spec{{\mathscr{S}}}

\def\Ideal{{\mathsf{Ideal}}}

%% Properties
\def\sfG{{\mathsf{G}}} % property G
\def\sfH{{\mathsf{H}}} % property H

%% Matrices
\def\sfA{{\mathsf{A}}} % the linear matrix
\def\GL{{\mathrm{GL}}} % the general linear group
\def\bfM{{\mathsf{M}}} % the matrix of change of variables
\def\rank{{\ensuremath{\mathsf{rank}}}} % rank function
\def\jac{{\ensuremath{\mathsf{jac}}}} % jacobian matrix
\def\minors{{\ensuremath{\mathsf{Minors}}}} % set of minors
\def\subm{{\ensuremath{\mathsf{SubMatrix}}}} % set of minors

%% Variables
\def\X{{X}} % variables of the space of the determinant
\def\Y{{Y}} % variables of the space of the incidence variety (kernel variables)
\def\Z{{Z}} % variables of the space of critical points (Lagrange multipliers)
\def\M{{M}}
\def\V{{V}}

%% Points
\def\x{{x}} % a generic first point
\def\y{{y}} % a generic second point
\def\z{{z}} % a generic third point
\def\v{{v}} % a generic third point
\def\fiber{{\alpha}} % a generic third point

%% Parameters
\def\fraka{{\mathfrak{a}}}
%\def\frakb{{\mathfrak{b}}}
%\def\frakc{{\mathfrak{c}}}
\def\fraku{{\mathfrak{u}}}
%\def\frakv{{\mathfrak{v}}}
\def\frakm{{\mathfrak{m}}}
\def\frakA{{\mathfrak{A}}}
%\def\frakB{{\mathfrak{B}}}
%\def\frakC{{\mathfrak{C}}}

\newcommand{\calV}{\mathcal{V}}

\def\Id{{\mathbb{I}}}

%% Algebraic sets
\def\setD{{\mathcal{D}}}
\def\setV{{\mathcal{V}}}
\def\setW{{\mathcal{W}}}
\def\spec{{\mathscr{S}}}

%% Zariski-open sets
\def\zarA{{\mathscr{A}}}
\def\zarfiber{{\mathscr{I}}}
\def\zarU{{\mathscr{U}}}
\def\zarV{{\mathscr{V}}}
\def\zarM{{\mathscr{M}}}
\def\zarU{{\mathscr{U}}}
\def\zarO{{\mathscr{O}}}

\def\softO{\ensuremath{{\mathcal{O}}{\,\tilde{ }\,}}}

\newcommand{\bX}{{x}}
\newcommand{\rY}{{y}}
\newcommand{\gZ}{{z}}

%% Messages
\def\xxx{{\textbf{(xxx)}}}
\def\why{\textbf{(why?)}}
\def\todo{\textbf{(todo!)}}
\def\check{\textbf{(check!)}}
\def\mymid{{\,\,:\,\,}}
\definecolor{mycolor}{rgb}{0.28,0.69,0.7}
\definecolor{mycolor1}{rgb}{0.7,0.75,0.71}
\definecolor{mycolor2}{rgb}{0.58,1,1}
\def\fall{\forall\,}
\def\MM{{\mathbb{M}}}
\def\SS{{\mathbb{S}}}
\def\sing{{\rm sing}\,}
\def\crit{{\rm crit}\,}
\def\reg{{\rm reg}\,}
\def\minors{{\rm minors}\,}
\def\spec{\mathscr{S}}
\def\cc{\mathcal{C}}

\usetikzlibrary{positioning}
\usetikzlibrary{fit}
\usetikzlibrary{backgrounds}
\usetikzlibrary{calc}
\usetikzlibrary{shapes}
%\usetikzlibrary{mindmap}
%\usetikzlibrary{decorations.text}

\tikzset{box1/.style={draw=black, thick, rectangle,rounded corners, minimum height=2cm, minimum width=2cm}}

\setbeamerfont{page number in head/foot}{size=\large}
\setbeamertemplate{footline}[frame number]
\setbeamercolor{page number in head/foot}{fg=black}

\begin{document}

%%%%%%%%%%%%%%%%%%%%%%%%%%%%%%%%%%%%%%%%%%%%%%%%%%%%%%%%%%%%%%%%%%%%%%%%%%%%%
\begin{frame}[noframenumbering]
\maketitle
  \begin{figure}[H]
    \centering
    \vspace{-0.2cm}
    \includegraphics[width=.5\linewidth]{../IMAGES/logoTU} \hfill\hbox{} 
%    \hspace{0.5cm}
 \hfill   \includegraphics[width=.4\linewidth]{../IMAGES/logomath}
  \end{figure}
\end{frame}

\begin{frame}
\frametitle{\vspace{-0.5cm}\bc \bf \Huge Spectrahedra and hyperbolicity cones \ec}
  \begin{figure}[H]
    \centering
    \vspace{-0.6cm}
    \includegraphics[width=.3\linewidth]{../IMAGES/intpoint}
    \hspace{1.5cm}
    \includegraphics[width=.4\linewidth]{../IMAGES/cayleyBig}\\
    \vspace{0.1cm}
    \includegraphics[width=.3\linewidth]{../IMAGES/honest_pillow}
    \hspace{0.3cm}
    \includegraphics[width=.3\linewidth]{../IMAGES/invQuartic}
    \hspace{0.1cm}
    \includegraphics[width=.3\linewidth]{../IMAGES/invQuintic}
  \end{figure}
\end{frame}




%%%%%%%%%%%%%%%%%%%%%%%%%%%%%%%%%%%%%%%%%%%%%%%%%%%%%%%%%%%%%%%%%%%%%%%%%%%%%
\begin{frame}
\frametitle{\vspace{-0.5cm}\bc \bf \Huge Positivity certificates of polynomials \ec}
Say I want to certify that
\[
f=U_1^4+U_1U_2^3+U_2^4-3U_1^2U_2U_3-4U_1U_2^2U_3+2U_1^2U_3^2+U_1U_3^3+U_2U_3^3+U_3^4
\]
is globally nonnegative ($f(u) \geq 0$ for all $u \in \RR^3$). \,\,\,\, 
%\pause
{\it E.g.}, is it a {\color{orange} \bf Sum Of Squares}?
%\pause
\vspace{0.2cm}
\bc Does there exist a PSD Gram matrix for $f$? \ec
\vspace{0.2cm}

\begin{tcolorbox}[colback=green!5,colframe=gray!40!black]
\setbeamercolor{mybox}{bg=orange}
\begin{beamercolorbox}[sep=1mm]{mybox}
$\text{With } \langle v \rangle_\RR = \RR[U]_2 \,\,\,\,\,\,\,\,\,\,\, \mathscr{S} = \{X \in \SS^6(\RR) \mymid f=v'Xv, X \succeq 0\} \neq \emptyset\,?$
\end{beamercolorbox}
\end{tcolorbox}
\vspace{0.4cm}
%\pause
\red{Bad news\,/\,Interesting facts} about this example:
\begin{itemize}
\item $f \in \sum (\RR[U]_2)^2$ but $f\notin\sum (\QQ[U]_2)^2$
\item Hence $\mathscr{S}$ is degenerate and $\mathscr{S} \cap \SS^6(\QQ) = \emptyset$
\item $f$ is never a sum of 3 non-zero squares (but it is of $2$)
\end{itemize}


\end{frame}


%%%%%%%%%%%%%%%%%%%%%%%%%%%%%%%%%%%%%%%%%%%%%%%%%%%%%%%%%%%%%%%%%%%%%%%%%%%%%
\begin{frame}
\frametitle{\vspace{-0.5cm}\bc \bf \Huge Linear matrix inequalities \ec}

%Given $C,M_1,\ldots,M_n \in \SS_m(\RR)$ and $b \in \RR^k$:
%  \begin{columns}[t,onlytextwidth]
%    \begin{column}{0.4\textwidth}
%      \bc{\bf Primal SDP}\ec
%      \begin{align*}
%        p^* = \inf_{X \in \mathbb{S}_m} & \,\,\, (C,X):=Trace(CX) \\
%         s.t.         & \,\,\, {\color{red} X \succeq 0} \\
%        & {\color{red} (M_j,X)=b_j, j=1, \ldots, k}
%      \end{align*}
%    \end{column}
%    \begin{column}{0.6\textwidth}
%      \bc{\bf Dual SDP}\ec
%      \begin{align*}
%        d^* = \sup_{s \in \RR^k} & \,\,\, b's \\
%        s.t.       & \,\,\, {\color{red} M(s) \succeq 0} \\
%         & {\color{red} M(s)=C-s_1M_1-\cdots-s_kM_k} 
%      \end{align*}
%    \end{column}
%  \end{columns}
%\pause
\vspace{-0.5cm}
\bc $A_0,A_1,\ldots,A_n$ are $m \times m$ real symmetric matrices \ec

\begin{tcolorbox}[colback=green!5,colframe=blue!40!black]
  \bc Spectrahedron: {\color{red} $\spec = \{x \in \RR^n \mymid A(x)=A_0+x_1A_1+\cdots+x_nA_n \succeq 0\}$} \ec
\end{tcolorbox}
It is {\bf basic semi-algebraic} since, if
$$\det(A(x)+t I_m) = f_{m}(x)+f_{m-1}(x)t+\cdots+f_{1}(x)t^{m-1}+t^m$$
then $\spec=\{x \in \RR^n \mymid f_i(x) \geq 0, i=1,\ldots,m\}$. $A(x) \succeq 0$ is called an {\bf LMI}.

\pause

\begin{columns}[t,onlytextwidth]
  \begin{column}{0.45\textwidth}
    \begin{figure}[H]
      \centering
      \includegraphics[width=.4\linewidth]{../IMAGES/cayleycubic}\\
      \includegraphics[width=.6\linewidth]{../IMAGES/cut}
      \caption{The Cayley spectrahedron and a maximum cut}
    \end{figure}
  \end{column}
  \begin{column}{0.6\textwidth}
    \bc {\bf SDP} : linear optimization over $\spec$ ({\it i.e.} over {\bf LMI}) \ec
    %      The {\bf Cayley spectrahedron}:
    \vspace{0cm}
    \[
    \spec=\left\{
    (x_1,x_2,x_3) \in \RR^3 \mymid
    \left(
    \begin{array}{ccc}
      1 & x_1 & x_2 \\
      x_1 & 1 & x_3 \\
      x_2 & x_3 & 1 \\
    \end{array}
    \right)\succeq 0
    \right\}
    \]
Applications in
\begin{itemize}
\item
control theory (Lyapunov stability ...)
\item
combinatorics (relaxations of NP-hard pr.)
\item
non-convex poly. optimization
\end{itemize}
%    \bc {\bf SDP relaxation of MAXCUT} \ec
%    \boxed{
%      \begin{array}{c}
%        {\color{red} \bf Quadratic\,\,opt.} \\
%        {\color{red} \bf problem\,\,(hard)}
%      \end{array}
%    }
%    $\longrightarrow$
%    \boxed{
%      \begin{array}{c}
%        {\color{blue} \bf Semidefinite} \\
%        {\color{blue} \bf program\,\,(easy)}
%      \end{array}
%    }
  \end{column}
\end{columns}

\end{frame}


%%%%%%%%%%%%%%%%%%%%%%%%%%%%%%%%%%%%%%%%%%%%%%%%%%%%%%%%%%%%%%%%%%%%%%%%%%%%%
\begin{frame}
\frametitle{\vspace{-0.5cm}\bc \bf \Huge Multivariate polynomial optimization \ec}

\begin{tcolorbox}[colback=green!50,colframe=blue!30!black,title=\bc {\bf The general problem} \ec]
  \bc $f^*=\displaystyle\inf_{x \in {\color{red} \bf S}} f(x)$, \,\,\, with \,\,\,{\color{red} ${\color{red} \bf S} = \{x \in \RR^n \mymid {\bf g_1}(x) \geq 0, \ldots, {\bf g_s}(x) \geq 0\}$} \ec
\end{tcolorbox}

\begin{itemize}
\item[$\bullet$] If ${\color{red} \bf S}=\RR^n \longrightarrow f^*=\displaystyle\sup_{\lambda \in \RR} \lambda$ \,\,
  s.t. \,\, $f(x)-\lambda \geq 0 \,\,\,\,\forall\,\,x\in \RR^n$ \,\, {\bf [NP-hard]}
\item[] Write $f-\lambda$ as a sum of squares? $f-\lambda= f_1^2+\cdots+f_p^2$ \,\,\, {\bf \color{blue} [SDP]}
\item[] Hilbert (1888): Positivity $\equiv$ SOS iff ${\color{olive} n=1}$ or ${\color{olive} d=2}$ or ${\color{olive} (n,d)=(2,4)}$
\pause
\vspace{0.2cm}
\item[$\bullet$] When ${\color{red} \bf S}$ is any set : {\bf Positivestellensatz}
\bc  (Schm\"udgen) \,\,${\color{red} \bf S}$ \,compact\,\,:\,\,
$f-y = \displaystyle\sum_{\nu \in \{0,1\}^s}\sigma_{\nu}\prod_j{\bf \color{red} g_j}^{\nu_j}$ \,\, with $\sigma_i$ SOS  \,\, {\bf \color{blue} [SDP]} \ec
\bc  (Putinar) \, \,\,\,\,${\color{red} \bf S}$ archimedean\,\,:\,\,\,
$f-y = \sigma_0 + \displaystyle\sum_i\sigma_i{\bf \color{red} g_i}$ \,\,\,\,\,\,\, with $\sigma_i$ SOS  \,\, {\bf \color{blue} [SDP]} \ec
\pause
\vspace{0.2cm}
\item[$\bullet$] Lasserre ('01), Parrilo ('03) \ldots Moments vs SOS approach also gives {\bf \color{blue} [SDP]}
\end{itemize}

\end{frame}


%%%%%%%%%%%%%%%%%%%%%%%%%%%%%%%%%%%%%%%%%%%%%%%%%%%%%%%%%%%%%%%%%%%%%%%%%%%%%
\begin{frame}
\frametitle{\vspace{-0.5cm}\bc \bf \Huge State of the art \ec}

\vspace{-0.3cm}

\setbeamercolor{mybox}{bg=gray}
\begin{beamercolorbox}[sep=1mm]{mybox}
${\bf Fixed\hspace{-0.1cm}-\hspace{-0.1cm}Precision \,\, Polynomial\hspace{-0.1cm}-\hspace{-0.1cm}Time \,\, Algorithms\,\, for\,\, SDP}$
\end{beamercolorbox}
{\cit Shor, Nemirovsky, Yudin ('60s-'70s), \ldots, }{\cit Khachiyan}\\
Ellipsoid algorithm is $Poly(m,n,\log(1/\epsilon))$ : {\color{red} inefficient in practice} \\
{\cit Nesterov, Nemirovsky ('90s)}\\
Interior point algorithms : {\color{blue} efficient in practice}

{\cit Ramana, Khachiyan, Porkolab ('90s)}\\
Complexity bounds for LMI : {\color{blue} can't be used in practice}

\pause

\vspace{0.1cm}

\setbeamercolor{mybox}{bg=gray}
\begin{beamercolorbox}[sep=1mm]{mybox}
${\bf (Rational)\,points\,in\,(convex)\,Semialgebraic\,Sets}$
\end{beamercolorbox}
{\cit Peyrl, Parrilo (2007), Kaltofen, Li, Yang, Zhi (2008)}\\
Floating-point-SOS $\rightarrow$ Rational approximation $\rightarrow$ Projection on the feasible set \\
{\cit Guo, Safey El Din, Zhi (2013)} \\
Bit complexity bounds for emptiness of $\{x \in \QQ^n \mymid \red{A}(x) \succeq 0\}$\\

\pause

\vspace{0.1cm}

\setbeamercolor{mybox}{bg=gray}
\begin{beamercolorbox}[sep=1mm]{mybox}
${\bf Exact\,\,algorithms\,\,for\,\,LMI/SDP}$
\end{beamercolorbox}
{\cit Henrion, N., Safey El Din (2015-2016)} \\
New upper bounds for emptiness of $\{x \in \RR^n \mymid \red{A}(x) \succeq 0\}$ \\
{\cit N. (2016)} \\
Rank-constrained SDP \hfill (poly-time if \, $n$ \, {\bf or} \, $m=\text{size}({A})$ \, is fixed)

%Open questions:
%\begin{itemize}
%\item
%How if $\{x \in \QQ^n \mymid \red{A}(x) \succeq 0\}=\emptyset$?
%\item
%Additional rank constraints?
%\end{itemize}


\end{frame}




%%%%%%%%%%%%%%%%%%%%%%%%%%%%%%%%%%%%%%%%%%%%%%%%%%%%%%%%%%%%%%%%%%%%%%%%%%%%%
\begin{frame}
\frametitle{\vspace{-0.5cm}\bc \bf \Huge Why algebraic methods? \ec}
\vspace{-0.8cm}
%\,\,\,\,\,\,\,\,\,\,\,\,{\cit Polyhedra and spectrahedra}
%
%      \begin{itemize}
%        \item {\color{red} \bf convex closed} sets
%        \item basic {\color{red} \bf semialgebraic:} $\spec = \{x \in \RR^n \mymid f_1(x) \geq 0, \ldots, f_s(x) \geq 0\}$
%        \item {\color{red} \bf exposed} faces: ${\cal F} = {\cal H} \cap \spec$
%        \item for $x \in \partial \spec \longrightarrow \det\,A(x)=0$
%        \item {\cit Nice property}: if $\cc \subset \calD_{\minrank}$ and $\cc \cap \spec \neq \emptyset$ then $\cc \subset \spec$
%      \end{itemize}
%          \pause
\bc 1. It is {\color{red} \bf Hard} to compute low-rank solutions to LMI/SDP \ec
\vspace{-0.2cm}
  \begin{columns}[t]
    \begin{column}[T]{5cm}
        \vspace{0.3cm}
      \begin{figure}[H]
        \centering
        \includegraphics[width=.35\linewidth]{../IMAGES/lp}
        \hspace{0.6cm}
        \includegraphics[width=.35\linewidth]{../IMAGES/honest_pillow}
        \label{fig:pillow}
        \caption{``{\bf Low-rank}'' points : they minimize a cone of linear forms}
      \end{figure}
    \end{column}
    \begin{column}[T]{5.2cm}
\vspace{0.3cm}
\begin{figure}[H]
  \centering
  \includegraphics[width=.5\linewidth]{../IMAGES/spectraExample-crop}
  \caption{{\sc sedumi} returns a floating point approximation of $(0,0)$ when maximizing $x_2$}
\end{figure}
\vspace{-0.2cm}
%      $$
%     \boxed{
%        \begin{array}{rl}
%          \inf &  (C,X) \\
%          s.t. &  X \succeq 0, \,\, (M_j,X)=b_j, j\leq k
%        \end{array}
%      }
%      $$
%      \begin{itemize}
%      \item {\color{red} \bf Hard} to find low-rank solutions to SDP
%      \item The set $\spec$ can have {\color{red} \bf empty interior}
%      \end{itemize}
    \end{column}
  \end{columns}

\bc 2. The interior of $\spec$ can be {\bf empty} $\longrightarrow$ {\bf \color{red} Interior point algorithms could fail} \ec

\[
      \begin{bmatrix}
        0 & x_1 & \frac{1}{2}(1-x_4) \\
        x_1 & x_2 & x_3 \\
        \frac{1}{2}(1-x_4) & x_3 & x_4
      \end{bmatrix} \succeq 0
\]

%  \begin{columns}[t,onlytextwidth]
%    \begin{column}{0.4\textwidth}
%\,\,\,\,\,\,\,\,\inf \,\,\,\,\,\,\,\,\,\,\,\, x_2\,\,$ : \\
%    \end{column}
%   \begin{column}{0.6\textwidth}
%      $
%      \begin{bmatrix}
%        0 & x_1 & \frac{1}{2}(1-x_4) \\
%        x_1 & x_2 & x_3 \\
%        \frac{1}{2}(1-x_4) & x_3 & x_4
%      \end{bmatrix} \succeq 0
%      $
%    \end{column}
%  \end{columns}


\end{frame}





%%%%%%%%%%%%%%%%%%%%%%%%%%%%%%%%%%%%%%%%%%%%%%%%%%%%%%%%%%%%%%%%%%%%%%%%%%%%%
\begin{frame}[noframenumbering]
\frametitle{\vspace{-0.5cm}\bc \bf \Huge Why algebraic methods? \ec}

\vspace{-1cm}

\begin{columns}[t,onlytextwidth]
  \begin{column}{0.5\textwidth}
\begin{figure}[H]
  \centering
  \includegraphics[width=.5\linewidth]{../IMAGES/spectraExample-crop}
\end{figure}
  \end{column}
  \begin{column}{0.5\textwidth}
\hbox{}\vspace{0.5cm}
\[
      \begin{bmatrix}
        0 & x_1 & \frac{1}{2}(1-x_4) \\
        x_1 & x_2 & x_3 \\
        \frac{1}{2}(1-x_4) & x_3 & x_4
      \end{bmatrix} \succeq 0
\]
  \end{column}
\end{columns}

\vspace{0.4cm}
Main motivations for the design of {\bf exact algorithms}:

\begin{enumerate}
\item Certification of output information (rank, algebraic degree...)
\item Manage {\bf algebraic constraints} such as rank defects?
\item Handle {\bf degenerate} non-full-dimensional examples?
\end{enumerate}

\pause

\begin{tcolorbox}[colback=blue!45,colframe=blue!40!white,title=\bf \color{black} Exact representation of a solution $x^*$ to $A(x) \succeq 0$ (RUR)]
\vspace{-0.3cm}
\[
(q,q_0,q_1,\ldots,q_{n}) \subset \QQ[t] \,\,\rightarrow \,\,x^* \in \left\{\left(\frac{q_1(t)}{q_0(t)}, \ldots, \frac{q_n(t)}{q_0(t)}\right) \mymid q(t)=0\right\}
\]
\end{tcolorbox}

\begin{itemize}
\item[]
{\cit Rouillier} [AAECC 1999]
\item[]
{\cit Giusti, Lecerf, Salvy} [JoC 2001] :: Geometric Resolution
\item[]
{\cit Jeronimo, Matera, Solernò, Waissbein} [FoCM 2009] :: Symbolic Homotopy
\end{itemize}

\end{frame}


%%%%%%%%%%%%%%%%%%%%%%%%%%%%%%%%%%%%%%%%%%%%%%%%%%%%%%%%%%%%%%%%%%%%%%%%%%%%%
\begin{frame}
\frametitle{\vspace{-0.5cm}\bc \bf \Huge The geometry of low-rank matrices \ec}

%For any $A(x)$ (not necessarily symmetric) we define
%\begin{itemize}
%\item[] {\it algebraic sets}\,\, $\calD_r = \{x \in \CC^n \mymid \rank\,A(x) \leq r\}, \, r=1,\ldots,m-1$
%\item For $A(x)$ symmetric, and $\spec \neq \emptyset$: $\minrank = \min \{ \rank\,A(x) \mid {x \in \spec} \}$
%\end{itemize}

Define:
\begin{itemize}
\item[] For any $A(x)$ (not nec. symmetric): ${\color{red} \calD_r = \{x \in \CC^n \mymid \rank\,A(x) \leq r\}}$
\item[] For $A(x)$ symmetric, and $\spec \neq \emptyset$: ${\color{red} \minrank = \min \{ \rank\,A(x) \mid {x \in \spec} \}}$
\end{itemize}

So one has a nested sequence : $\calD_0\cap\RR^n \subset \cdots \subset \calD_{m-1}\cap\RR^n$

\vspace{0.3cm}

  \begin{tcolorbox}[colback=green!5,colframe=blue!40!black,title=Smallest Rank Property (Henrion\,-\,N.\,-\,Safey El Din)]
    $A(x)$ symmetric, and $\spec \neq \emptyset$. Let $\cc$ be a conn. comp. of $\setD_{\minrank} \cap \RR^n$
    s.t. $\cc \cap \spec \neq \emptyset$. \,\,\,\,{\bf Then} $\cc \subset \spec$. \,\,\,\,{\bf In particular}
    \,$\cc \subset \setD_{\minrank} \setminus \setD_{\minrank-1}$.
  \end{tcolorbox}

\vspace{0.2cm}

\pause

$\rightarrow$ Computation of {\bf one point per connected component} of algebraic sets :

\begin{itemize}
\item[]
{\small \cit Safey El Din, Schost} [ISSAC 2003, DCG 2004]
\item[]
{\small \cit Bank, Giusti, Heintz, Safey El Din, Schost} [AAECC 2010]
\end{itemize}

\pause

$\rightarrow$ Exploiting the {\bf determinantal structure} for polynomial system solving :

\begin{itemize}
\item[]
{\small \cit Bank, Giusti, Heintz, Lecerf, Matera, Solernó} [Found. Comput. Math. (2015)]
\item[]
{\small \cit Henrion, N., Safey El Din} [JSC 2015, ISSAC 2015]
\end{itemize}


%  $$
%  \boxed{
%    \begin{array}{c}
%    {\color{black} \bf Either} \\
%    \,\,\text{the set $\spec$ is empty}
%    \end{array}
%  }
%  \hspace{1cm}
%  \boxed{
%    \begin{array}{c}
%    {\color{black} \bf Or} \\
%    \text{$\minrank$ is well-defined} \\
%    \text{$\exists\,\cc \subset \calD_{\minrank} \mymid \cc \subset \spec$}  
%    \end{array}
%  }
%  $$

\end{frame}




%%%%%%%%%%%%%%%%%%%%%%%%%%%%%%%%%%%%%%%%%%%%%%%%%%%%%%%%%%%%%%%%%%%%%%%%%%%%%
\begin{frame}
\frametitle{\vspace{-0.5cm}\bc \bf \Huge Incidence varieties and critical points \ec}

\vspace{0.3cm}
  \begin{columns}[t,onlytextwidth]
    \begin{column}{0.6\textwidth}
      {\bf\underline{1st ste}p} ${\it Lifting\,of\,the\,determinantal\,variety\,\calD_r}$:
      $$
      \sfA( {x} ) \, Y({\color{red} y})
      = \sfA( {x} ) \, 
      \left[
        \begin{array}{ccc}
          {\color{red} \y_{1,1}} & \ldots & {\color{red} \y_{1,m-r}} \\ 
          \vdots    &        & \vdots      \\ 
          {\color{red} \y_{m,1}} & \ldots & {\color{red} \y_{m,m-r}} \\ 
        \end{array}
        \right]
      = 0.
      $$
%      $$
%      U \, {{Y({\color{red} y})}} = I_{m-r}
%      $$
    \end{column}
    \begin{column}{0.5\textwidth}
\vspace{-1cm}
      \begin{figure}[h]
        \hspace{5cm}
        \includegraphics[scale=0.55]{../IMAGES/component1}
      \end{figure}
    \end{column}
  \end{columns}
  \begin{tcolorbox}[colback=green!5,colframe=blue!40!black]
    \bc If $A$ is generic, the lifted algebraic set ${\cal V}_r$ is {\bf smooth} and {\bf equidimensional} \ec
  \end{tcolorbox}
\vspace{0.2cm}

  {\bf\underline{2nd ste}p} ${\it Compute\,critical\,points}$ of the map $\pi(x,{\color{red} y}) = a_1x_1+\cdots+a_nx_n$ on ${\cal V}_r$:
  \begin{tcolorbox}[colback=green!5,colframe=blue!40!black]
    \bc When $A$ and $a_1\,..\,a_n$ are generic, there are {\bf finitely many} critical points. \ec
  \end{tcolorbox}
\vspace{0.2cm}

  {\bf\underline{3rd ste}p} ${\it Intersect\,with\,any\,fiber\,of\,\pi\,and\,call\,point\,2}$
\end{frame}


%%%%%%%%%%%%%%%%%%%%%%%%%%%%%%%%%%%%%%%%%%%%%%%%%%%%%%%%%%%%%%%%%%%%%%%%%%%%%
\begin{frame}[noframenumbering]
\frametitle{\vspace{-0.5cm}\bc \bf \Huge Incidence varieties and critical points \ec}

\vspace{0.3cm}
  \begin{columns}[t,onlytextwidth]
    \begin{column}{0.6\textwidth}
      {\bf \underline{1st ste}p} ${\it Lifting\,of\,the\,determinantal\,variety\,\calD_r}$:
      $$
      \sfA( {x} ) \, Y({\color{red} y})
      = \sfA( {x} ) \, 
      \left[
        \begin{array}{ccc}
          {\color{red} \y_{1,1}} & \ldots & {\color{red} \y_{1,m-r}} \\ 
          \vdots    &        & \vdots      \\ 
          {\color{red} \y_{m,1}} & \ldots & {\color{red} \y_{m,m-r}} \\ 
        \end{array}
        \right]
      = 0.
      $$
%      $$
%      U \, {{Y({\color{red} y})}} = I_{m-r}
%      $$
    \end{column}
    \begin{column}{0.5\textwidth}
\vspace{-1cm}
      \begin{figure}[h]
        \hspace{5cm}
        \includegraphics[scale=0.55]{../IMAGES/component2}
      \end{figure}
    \end{column}
  \end{columns}
  \begin{tcolorbox}[colback=green!5,colframe=blue!40!black]
    \bc If $A$ is generic, the lifted algebraic set ${\cal V}_r$ is {\bf smooth} and {\bf equidimensional} \ec
  \end{tcolorbox}
\vspace{0.2cm}

  {\bf \underline{2nd ste}p} ${\it Compute\,critical\,points}$ of the map $\pi(x,{\color{red} y}) = a_1x_1+\cdots+a_nx_n$ on ${\cal V}_r$:
  \begin{tcolorbox}[colback=green!5,colframe=blue!40!black]
    \bc When $A$ and $a_1\,..\,a_n$ are generic, there are {\bf finitely many} critical points. \ec
  \end{tcolorbox}
\vspace{0.2cm}

  {\bf\underline{3rd ste}p} ${\it Intersect\,with\,any\,fiber\,of\,\pi\,and\,call\,point\,2}$
\end{frame}





%%%%%%%%%%%%%%%%%%%%%%%%%%%%%%%%%%%%%%%%%%%%%%%%%%%%%%%%%%%%%%%%%%%%%%%%%%%%%
\begin{frame}
\frametitle{\vspace{-0.5cm}\bc \bf \Huge {\sc spectra}: a library for linear matrix inequalities \ec}

\vspace{-0.4cm}

%What is {\sc spectra}? \\

\begin{itemize}
%\item[] {\bf S}emidefinite {\bf P}rogramming and {\bf E}xact {\bf C}omputation {\bf T}owards {\bf R}eal {\bf A}lgebra
\item[] A {\sc maple} library, freely distributed %%%%, $\approx 3000$ lines of code 
\item[] Depends on {\sc fgb} for computations with Gr\"obner bases ({J-C. Faugère})
\item[] Addressed to researchers in {Optimization, Convex alg. geom., Symb. comp.}
\end{itemize}

\vspace{-0.2cm}

\begin{columns}[t,onlytextwidth]
  \begin{column}{0.5\textwidth}
\vspace{-0.4cm}   
    \begin{table}[!ht]
      \centering
          {\small
            \begin{tabular}{|l|rrr|}
              \hline
              \rowcolor{mycolor} $(m,r,n)$ & {\sc raglib} & {\sc spectra} & {\sf deg} \\
              \hline
              \hline
              \rowcolor{mycolor2} $(3,2,8)$ & 109      & 18    & 39 \\
              $(3,2,9)$ & 230      & 20    & 39 \\
              \rowcolor{mycolor2} $(4,2,5)$ & 12.2     & 26    & 100 \\
              $(4,2,6)$ & $\infty$ & 593    & 276  \\
              \rowcolor{mycolor2} $(4,2,7)$ & $\infty$ & 6684   & 532 \\
              $(4,2,8)$ & $\infty$ & 42868  & 818  \\
              \rowcolor{mycolor2} $(4,2,9)$ & $\infty$ & 120801 & 1074 \\
              $(4,3,10)$ & $\infty$ & 303  & 284 \\
              \rowcolor{mycolor2} $(4,3,11)$ & $\infty$ & 377  & 284  \\
%              $(5,2,7)$ & 4.2   & 0.7  &  0  \\
%              \rowcolor{mycolor2} $(5,2,8)$ &  8  & 0.7  & 0 \\
              $(5,2,9)$ & $\infty$   & 903 & 175 \\
              \rowcolor{mycolor2} $(6,5,4)$ & $\infty$ & 8643  & 726 \\
              \hline
            \end{tabular}
          }
    \end{table}
    
  \end{column}
  \begin{column}{0.55\textwidth}
    \vspace{0.3cm}
    \begin{itemize}
    \item[-] {\sc raglib} = {\cit Safey El Din}'s real algebraic geometry library
      \vspace{0.1cm}
    \item[-] {\sf deg} = degree of Rational Parametrization
      \vspace{0.1cm}
    \item[-] Time in seconds
      \vspace{0.1cm}
    \item[-] $\infty =$ more than 2 days
    \end{itemize}
  \end{column}
\end{columns}
\vspace{-0.2cm}
\bc

\vfill
\begin{tcolorbox}[colback=blue!45,colframe=blue!40!white]
Web page : \,\,\,\,{\tt \bf \small www.mathematik.tu-dortmund.de/sites/simone-naldi/software}
\end{tcolorbox}
\ec

\end{frame}

%%%%%%%%%%%%%%%%%%%%%%%%%%%%%%%%%%%%%%%%%%%%%%%%%%%%%%%%%%%%%%%%%%%%%%%%%%%%%
\begin{frame}
\frametitle{\vspace{-0.5cm}\bc \bf \Huge {\sc spectra}: a library for linear matrix inequalities \ec}

In a Maple worksheet type :

\vspace{0.2cm}

{\tt
> with(SPECTRA);

\hspace{0.3cm} [SolveLMI]
}

\[
A(x) = 
\left[
\begin{array}{ccc}
1+x_1 & x_2 & 0 \\
x_2 & 1-x_1 & 0 \\
0 & 0 & x_1
\end{array}
\right]
\succeq 0
\begin{array}{c}
\text{ This LMI defines the {\bf half-disk} } \\
\spec =
\{
(x_1,x_2) : x_1^2+x_2^2 \leq 1, x_1 \geq 0
\}
\end{array}
\]
{\tt
%> A := Matrix([[1+x1,x2,0], [x2,1-x1,0], [0,0,x1]]):
%
> SolveLMI(A,\{rnk\});             

 \hspace{0.3cm} [[x1 = [0, 0], x2 = [1, 1]], rnk = 1]
}

\pause

\[
A(x) = 
\left[
\begin{array}{ccc}
1+x_1 & x_2 & 0 \\
x_2 & 1-x_1 & 0 \\
0 & 0 & {\color{red} \bf x_1-1}
\end{array}
\right]
\succeq 0
\begin{array}{c}
\text{ This LMI defines {\bf one single point} } \\
\spec =
\{
(1,0)
\}
\end{array}
\]

{\tt
> SolveLMI(A,\{rnk,deg\});

 \hspace{0.3cm}  [[x1 = [1, 1], x2 = [0, 0]], rnk = 1, deg = 1]
}



%\vfill
%\begin{tcolorbox}[colback=blue!45,colframe=blue!40!white]
%Web page : \,\,\,\,{\tt \bf \small www.mathematik.tu-dortmund.de/sites/simone-naldi/software}
%\end{tcolorbox}
%\ec

\end{frame}


%%%%%%%%%%%%%%%%%%%%%%%%%%%%%%%%%%%%%%%%%%%%%%%%%%%%%%%%%%%%%%%%%%%%%%%%%%%%%
\begin{frame}
\frametitle{\vspace{-0.5cm}\bc \bf \Huge {\sc spectra}: a library for linear matrix inequalities \ec}


\[
A(x) = 
\left[
\begin{array}{ccc}
1+x_1 & x_2 & 0 \\
x_2 & 1-x_1 & 0 \\
0 & 0 & {\color{red} \bf x_1-1-10^{-20}}
\end{array}
\right]
\succeq 0
\begin{array}{c}
\,\,\,\,\, \text{ This LMI defines } \spec = \emptyset
\end{array}
\]
{\tt
%> A := Matrix([[1+x1,x2,0], [x2,1-x1,0], [0,0,x1]]):
%
> SolveLMI(A);

 \hspace{0.3cm} [ ]
}

\pause

\[
A(x) = 
\left[
\begin{array}{ccc}
1+x_1 & x_2 & 0 \\
x_2 & 1-x_1 & 0 \\
0 & 0 & {\color{red} \bf x_1-1+10^{-20}}
\end{array}
\right]
\succeq 0
\begin{array}{c}
\,\,\,\, \text{ This LMI defines} \\ 
\,\,\,\, \textbf{a tiny half-disk}  \\
\end{array}
\]

{\tt
> SolveLMI(A);

 \hspace{0.3cm} [[x1 = [36893488147418995335 / 36893488147419103232,

 \hspace{1.8cm}  4611686018427401391 / 4611686018427387904],

 \hspace{0.7cm}  x2 = [-350142318592414077 / 

\hfill 2475880078570760549798248448,

 \hspace{1.8cm} -2801138548739304423 /

\hfill 19807040628566084398385987584]]
}


\end{frame}



%%%%%%%%%%%%%%%%%%%%%%%%%%%%%%%%%%%%%%%%%%%%%%%%%%%%%%%%%%%%%%%%%%%%%%%%%%%%%
\begin{frame}
\frametitle{\vspace{-0.5cm}\bc \bf \Huge {\sc spectra}: a library for linear matrix inequalities \ec}

A {\bf quartic plane curve} $\{x \in \RR^2 : \det A(x_1,x_2) = 0\}$ \\

\vspace{0.2cm}

The {\bf convex set} is $\{x \in \RR^2 : A(x_1,x_2) \succeq 0\}$

  \begin{figure}[H]
    \centering
    \includegraphics[width=.45\linewidth]{../IMAGES/quartic-crop}
  \end{figure}

Iterate the following command to compute many sample points :

{\tt
> SolveLMI(A,\{\},[3]);
}



\end{frame}




\begin{frame}
\frametitle{\vspace{-0.5cm}\bc \bf \Huge {\sc spectra}: a library for linear matrix inequalities \ec}

\[
f=U_1^4+U_1U_2^3+U_2^4-3U_1^2U_2U_3-4U_1U_2^2U_3+2U_1^2U_3^2+U_1U_3^3+U_2U_3^3+U_3^4
\]

Its Gram matrix satisfies $A(x) \succeq 0, f = v^T A(x) v$, and equals :

\vspace{0.1cm}

{\tt
> A := Matrix([[1,0,x1,0,-3/2-x2,x3], [0,-2*x1,1/2,x2,-2-x4,-x5], [x1,1/2,1,x4,0,x6],
               [0,x2,x4,-2*x3+2,x5,1/2], [-3/2-x2,-2-x4,0,x5,-2*x6,1/2], [x3,-x5,x6,1/2,1/2,1]]):
}

\vspace{0.2cm}

A minimal certificate is given by :

\vspace{0.1cm}

{\tt
> SolveLMI(A, \{rnk, deg, par\});

[[[x1 = [..], x2 = [..], x3 =  [..], x4 = [..], x5 = [..], x6 = [..]],
    rnk = 2, deg = 3,
    par = [8*z\^{}3+8*z+1, 24*z\^{}2-8, [16*z+3, -24*z\^{}2+8, 8*z\^{}2+6*z+8, -16*z\^{}2+6*z+16, -16*z-3, 16*z+3]]]
}

\end{frame}






\begin{frame}
\frametitle{\vspace{-0.5cm}\bc \bf \Huge Performances \ec}

\begin{figure}[!ht]
\centering
\begin{tikzpicture}
\begin{loglogaxis} [xlabel = number of variables, ylabel = time (seconds)]
\addplot [ultra thick, red]
file {m2r1.txt};
\addplot [ultra thick, blue]
file {m3r2.txt};
\end{loglogaxis}
\end{tikzpicture}
\caption{Random LMI of size $m={\bf \color{red} 2}$ and
$m = {\bf \color{blue} 3}$ and rank $r=m-1$, as a function of $n$.}
\label{fig:m2r1}
\end{figure}

\end{frame}


\begin{frame}
\frametitle{\vspace{-0.5cm}\bc \bf \Huge Exponential Bit-size \ec}

For a given $n \in {\mathbb N}$, consider the spectrahedron
\[
{\mathscr S}_n = \left\{x \in {\mathbb R}^n :
\left[\begin{array}{cc}
1 & 2 \\
2 & x_1
\end{array}\right] \succeq 0,
\left[\begin{array}{cc}
1 & x_1 \\
x_1 & x_2
\end{array}\right] \succeq 0,
\cdots, 
\left[\begin{array}{cc}
1 & x_{n-1} \\
x_{n-1} & x_n
\end{array}\right]\succeq 0
\right\}.
\]
For every $x^* \in {\mathscr S}_n$ it holds
$$
x^*_n \geq (x^*_{n-1})^2 \geq \cdots \geq (x^*_1)^{2^{n-1}} \geq 2^{2^n}.
$$

%%% \,\,\,\, {\tt Intel i7 / 2.5GHz / 16GB RAM}

\pause

\vspace{0.2cm}

I call {\sc spectra} with :

\vspace{0.2cm}

{\tt
> SolveLMI(A,\{\},[n]);

\hspace{0.5cm}       [[x1 = 4, x2 = 16, ..., x6 = 18446744073709551616]]
}

\vspace{0.2cm}

{\bf Average Timings} :
\begin{itemize}
\item[] $2n = 10 \rightarrow 15$ seconds
\item[] $2n = 12 \rightarrow 505$ seconds
\end{itemize}

This is due to the ``sparsity'' of this linear matrix.

\end{frame}



\begin{frame}
\frametitle{\vspace{-0.5cm}\bc \bf \Huge Hyperbolic polynomials \ec}

\begin{tcolorbox}[colback=white!50,colframe=blue!30!black,title=\bc {\bf Definition} \ec]
  \begin{columns}[t,onlytextwidth]
    \begin{column}{0.6\textwidth}
      $f \in \RR[x]_d$ is {\it hyperbolic w.r.t.} $e \in \RR^n$ if
      \begin{itemize}
      \item[(1)] $f(e) > 0$
      \item[(2)] $\forall\,a \in \RR^n \,\,\,\, t \mapsto f(t\,e-a)$ \\  has {\bf only real roots}
      \end{itemize}
          {\bf {\color{red} [LP] }  $f = \ell_1(x) \cdots \ell_d(x)$, $\,\,\,\,\ell_i \in \RR[x]_1$}\\
          {\bf {\color{red} [SDP] } Symmetric Determinant $f=\det(X)$}
    \end{column}  
    \begin{column}{0.4\textwidth}
      \begin{figure}[h]
        \includegraphics[scale=0.5]{../IMAGES/hyperbolic}
      \end{figure}
    \end{column}
  \end{columns}
\end{tcolorbox}

\pause

\vspace{0.2cm}

{\bf Hyperbolicity cone : } $\Lambda_+(f) = \{ a \in \RR^n : f(t\,e-a) = 0 \Rightarrow t \geq 0 \}$ {\bf \color{red} EX:} PSD cone

\vspace{0.2cm}

{\bf \color{red} Hyperbolic Programming [HP] : } Optimizing linear functions on $\Lambda_+$ (PDEs theory)

\vspace{0.2cm}

\pause

\begin{tcolorbox}[colback=green!5,colframe=orange!40!black]
  {\bf Generalized Lax conjecture:} \hfill Is every hyperbolicity cone a spectrahedron?
\end{tcolorbox}

\vspace{0.2cm}

{\cit N., Plaumann} [Preprint 2016] $\rightarrow$ Computer algebra techniques for HP

\end{frame}



%%%%%%%%%%%%%%%%%%%%%%%%%%%%%%%%%%%%%%%%%%%%%%%%%%%%%%%%%%%%%%%%%%%%%%%%%%%%%%
%\begin{frame}
%\frametitle{\vspace{-0.5cm}\bc \bf \Huge Spectrahedral shadows \ec}
%
%  \begin{columns}[t,onlytextwidth]
%    \begin{column}{0.5\textwidth}
%\begin{center}
%\vspace{-0.8cm}
%%{\bf \Large Spectrahedral shadows:}
%\end{center}
%\[
%S = \left\{ x \in \RR^n \mymid \exists\, u \in \RR^{d}, A(x, u) \succeq 0\right\},
%\]
%
%\begin{center}
%{\bf Projection on $\RR^n$ of the spectrahedron:}
%\end{center}
%
%$\spec = \{(x, u) \in \RR^n\times\RR^d \mymid A(x, u) \succeq 0\}$
%
%    \end{column}
%    \hspace{5pt}
%    \begin{column}{0.5\textwidth}
%\vspace{-1.4cm}
%  \begin{figure}[h]
%    \hspace{5cm}
%    \includegraphics[scale=0.17]{../IMAGES/proj}
%  \end{figure}
%
%    \end{column}
%  \end{columns}
%
%\vspace{0.1cm}
%
%  \begin{tcolorbox}[colback=green!5,colframe=orange!40!black]
%   {\bf Helton-Nie conjecture:} Are all convex semialgebraic sets spectrahedral shadows?
%   \hspace{3.5cm} {\cit Scheiderer}: true in dimension $n=2$
%  \end{tcolorbox}
%\vspace{0.3cm}
%
%{\bf Remark!} \,\,\, {Lin. Opt. on $S$} $\equiv$ {Lin. Opt. on $\spec$} $\equiv$ {\bf SDP}! \, {\color{red} \bf \Large \hfill (But $d \gg n^{\bigO(1)}$)}
%\vspace{0.2cm}
%
%  \begin{tcolorbox}[colback=green!5,colframe=green!40!black]
%   {\bf Perspective:} \hbox{}\hfill Parametric study of determinantal varieties \\
%   \hbox{}\hfill Solve linear optimization over $S$, given $\spec$ ?
%  \end{tcolorbox}
%
%\end{frame}


\begin{frame}
\frametitle{\vspace{-0.5cm}\bc \bf \Huge Perspectives \ec}

{\bf Structured problems from applied contexts}
\begin{itemize}
\item
  Lyapunov Stability :: Block-LMI conditions of type
  $$ 
  \begin{array}{rr}
    P \succ 0 \\
    PA+A^TP \prec 0
  \end{array}
  \Longleftrightarrow
  \begin{bmatrix}
    \begin{array}{cc}
      P & 0 \\
      0 & -PA-A^TP
    \end{array}
  \end{bmatrix}
  \succ 0
  $$
\end{itemize}

\vspace{0.4cm}

{\bf Towards an exact solver for SDP}

\begin{itemize}
\item
  No assumptions on input data using deformation/homotopy techniques
\end{itemize}

\vspace{0.4cm}

{\bf Exact algorithms for hyperbolic polynomials}

\begin{itemize}
\item
  Hyperbolicity test for $f \in \RR[x]_d$
\item
  Determinantal representations $\leftrightarrow$ cone of {\it interlacers} $I(f,e)$
\end{itemize}

%\vspace{0.2cm}
%
%{\bf Performance of the implementation}
%
%\begin{itemize}
%\item
%  Gr\"obner bases of multi-homogeneous Lagrange systems
%\item
%  Use/Implement routines for RUR (symbolic homotopy, geometric resolution ...)
%\end{itemize}



\end{frame}




\end{document}
